\chapter{Sinus, cosinus et oscillations}
\section*{Exercice 17.1 -- Angles et signes}
[Base] Calculer $\cos(5\pi/3)$, $\sin(-\pi/6)$ et résoudre $\cos x=1/2$ dans $[-\pi,\pi]$.

\section*{Exercice 17.2 -- Inéquation}
[Méthode] Résoudre $\cos x\ge-1/2$ dans $[-\pi,\pi]$.

\section*{Exercice 17.3 -- Composition}
[Base] Dériver $\sin(3x^2)$ et étudier $1+\cos(2x)$ sur $[0,\pi]$.

\section*{Exercice 17.4 -- Rectangle inscrit}
[Modélisation] Un rectangle est inscrit dans un cercle de rayon $R>0$. En paramétrant ses demi-côtés par $R\cos t,R\sin t$, $0<t<\pi/2$, maximiser son aire.

\section*{Exercice 17.5 -- Limite trigonométrique}
[Preuve] Montrer $\lim_{x\to0}\sin(5x)/x=5$ et $\lim_{x\to0}(1-\cos x)/x^2=1/2$.

\section*{Exercice 17.6 -- Tangente et domaine}
[Approfondissement] Définir $\tan x$, dériver cette fonction et expliquer pourquoi on ne peut pas conclure qu'elle est croissante sur tout son domaine.
